% Source: https://artofproblemsolving.com/community/c3h2239260_fisher_fishing
% Stable identity: aops-topic-2239260-q01
The probabilities of catching a fish on each day are
\[\text{Mon: }0.8\cdot0.8=\dfrac{16}{25},\quad\text{Tue: }0.8\cdot0.3=\dfrac{6}{25},\quad\text{Wed: }0,\quad\text{Thu--Sun: }0.3\cdot0.2=\dfrac{3}{50}\text{ each}.\]
Use the complement: no two consecutive days with a catch. Since there is never a catch on Wednesday, the days Monday--Tuesday and Thursday--Sunday can be treated separately.
\begin{itemize}
\item Monday and Tuesday are not both catch days with probability $1-\dfrac{16}{25}\cdot\dfrac{6}{25}=\dfrac{529}{625}$.
\item Thursday to Sunday: let $p=\dfrac{3}{50}$, $q=1-p=\dfrac{47}{50}$. With no two consecutive catches among $4$ days, there are no catches, one catch ($4$ ways) or two catches on non-consecutive days ($\{1,3\},\{1,4\},\{2,4\}$). The probability is
\[q^4+4pq^3+3p^2q^2=q^2(q+p)(q+3p)=q^2(1+2p)=\dfrac{2209}{2500}\cdot\dfrac{56}{50}=\dfrac{15463}{15625}.\]
\end{itemize}
So the probability of no two consecutive catch days is $\dfrac{529}{625}\cdot\dfrac{15463}{15625}=\dfrac{8179927}{9765625}$, and the required probability is
\[1-\dfrac{8179927}{9765625}=\boxed{\dfrac{1585698}{9765625}}\approx0.162.\]
